\section{Introduction}

Example 7.6 of the reader cites the experiment by Abram and Fieg, where the
adaptive controller is implemented on an Arduino Due \cite{ravazzoloMehrle2026,abramFieg}.
The present report realises the same self-tuning pole-placement idea on an
own test rig. The hardware consists of an ESP32, an L298N H-bridge, and a DC
motor with a quadrature encoder. The controller is not hand-coded on the
microcontroller; instead, the real-time code is generated automatically from
the Simulink model.

The structure follows the indirect adaptive-control scheme of Figure 7.1 in
the reader: the plant parameters are estimated recursively, the controller is
redesigned from these estimates, and the resulting incremental PI algorithm is
executed on the ESP32 input/output interface.
Figure~\ref{fig:self-tuning-regulator} shows the arrangement for the motor
rig. After calibration by the one-revolution encoder test, the first-order
self-tuning loop was run on the rig. The second-order controller was also run
on the rig, but is reported only as diagnostic evidence for why full-order
identification fails there; the working second-order design is a simulation
result.

\begin{figure}[htbp]
  \centering
  \resizebox{0.85\linewidth}{!}{%
  \begin{tikzpicture}[
      auto,
      >=Latex,
      node distance=12mm and 14mm,
      block/.style={draw, rectangle, minimum height=9mm, minimum width=27mm, align=center},
      smallblock/.style={draw, rectangle, minimum height=8mm, minimum width=24mm, align=center},
      sum/.style={draw, circle, inner sep=0pt, minimum size=5.5mm},
      signal/.style={->, thick}
    ]
    \node[sum] (sum) {};
    \node[left=10mm of sum] (r) {$r$};
    \node[block, right=13mm of sum] (pi) {incremental\\PI};
    \node[block, right=18mm of pi] (io) {ESP32\\I/O};
    \node[block, right=18mm of io] (plant) {DC motor\\and encoder};
    \node[right=22mm of plant] (y) {$y$};

    \node[smallblock, above=26mm of io] (pp) {pole\\placement};
    \node[smallblock, above=26mm of plant] (rls) {RLS\\estimator};

    \draw[signal] (r) -- (sum);
    \draw[signal] (sum) -- node {$e$} (pi);
    \draw[signal] (pi) -- node {$u$} (io);
    \draw[signal] (io) -- (plant);
    \draw[signal] (plant) -- (y);

    \coordinate (fbtap) at ($(plant.east)+(6mm,0)$);
    \coordinate (ytap)  at ($(plant.east)+(13mm,0)$);
    \coordinate (utap)  at ($(pi.east)+(5mm,0)$);
    \coordinate (uin)   at ($(rls.south)+(-6mm,0)$);
    \coordinate (yin)   at ($(rls.south)+(6mm,0)$);

    \draw[signal] (fbtap) |- ++(0,-13mm) -| node[pos=0.95, left] {$-$} (sum.south);

    \draw[signal] (utap) |- ($(uin)+(0,-11mm)$) -- (uin);
    \draw[signal] (ytap) |- ($(yin)+(0,-6mm)$) -- (yin);

    \fill (utap)  circle (0.6mm);
    \fill (ytap)  circle (0.6mm);
    \fill (fbtap) circle (0.6mm);

    \draw[signal] (rls) -- node {$\hat a_0,\hat b_0$} (pp);

    \coordinate (dcorner) at ($(pp.south)+(0,-5mm)$);
    \coordinate (dturn)   at (pi.north |- dcorner);
    \draw[signal] (pp.south) -- (dcorner)
      -- node[above] {$d_0,d_1$} (dturn) -- (pi.north);

    \node[above=4mm of pp] {controller design};
    \node[above=4mm of rls] {plant estimation};
  \end{tikzpicture}
  }
  \caption{Self-tuning regulator used for the ESP32 DC motor rig.}
  \label{fig:self-tuning-regulator}
\end{figure}
